\documentclass[11pt]{article}
\usepackage{explorer}
\newcommand{\ExplorerLanguage}{EN}
\newcommand{\ExplorerNumber}{QUADRATIC}
\newcommand{\ExplorerKind}{Present and explore}
\titlespacing*{\section}{0pt}{8pt}{4pt}
\setlength{\parskip}{5pt}\linespread{1.06}
\begin{document}
\ExplorerTitle{Upper high school}{Quadratic polynomials}
\ExplorerSection{Expanded form}
For $f(x)=ax^2+bx+c$, with $a,b,c\in\mathbb R$ and $a\ne0$, the domain is $\mathbb R$. The sign of $a$ determines the opening of the parabola.\par
\ExplorerSection{Vertex form}
$f(x)=a(x-\alpha)^2+\beta$, where $\alpha=-\dfrac{b}{2a}$ and $\beta=f(\alpha)=c-\dfrac{b^2}{4a}$. The vertex is $S(\alpha,\beta)$ and the axis of symmetry is $x=\alpha$. If $a>0$, $f$ decreases on $(-\infty,\alpha]$ and increases on $[\alpha,+\infty)$; $\beta$ is its minimum. If $a<0$, the directions reverse and $\beta$ is its maximum.\par
\ExplorerSection{Discriminant and roots}
$\Delta=b^2-4ac$ and $\beta=-\dfrac{\Delta}{4a}$. If $\Delta>0$, the two roots are $\dfrac{-b-\sqrt\Delta}{2a}$ and $\dfrac{-b+\sqrt\Delta}{2a}$: order them before naming them $x_1<x_2$ (the order of these two formulas depends on the sign of $a$). If $\Delta=0$, the double root is $x_0=\alpha$. If $\Delta<0$, there are no real roots.\par
\ExplorerSection{Factorisation and sign}
If $\Delta>0$, $f(x)=a(x-x_1)(x-x_2)$. Outside the roots the two factors have the same sign; between them they have opposite signs. Thus $f$ has the sign of $a$ outside, the opposite sign between, and is zero at the roots. If $\Delta=0$, $f(x)=a(x-x_0)^2$ has the sign of $a$ except at $x_0$. If $\Delta<0$, $f$ has strictly the sign of $a$ throughout $\mathbb R$ and cannot be factorised into real linear factors.\par
\ExplorerSection{Worked example}
$2x^2-4x-6=2(x-1)^2-8=2(x+1)(x-3)$. The vertex is $(1,-8)$, $\Delta=64$ and the roots are $-1$ and $3$. Hence $f(x)\le0$ is equivalent to $x\in[-1,3]$.\par
\begin{center}\begin{tikzpicture}[x=.65cm,y=.13cm]\draw[->,gray](-3,0)--(5,0)node[right]{$x$};\draw[->,gray](0,-10)--(0,12)node[above]{$y$};\draw[violet,thick,domain=-2.15:4.15,samples=100]plot(\x,{2*(\x-1)^2-8});\draw[mint,dashed](1,-9)--(1,11);\fill[mint](-1,0)circle(2pt)node[below]{$-1$};\fill[mint](3,0)circle(2pt)node[below]{$3$};\fill[violet](1,-8)circle(2pt)node[right]{$S(1,-8)$};\end{tikzpicture}\end{center}
\ExplorerSection{Investigate and justify}
Find the variation, roots and sign of $-x^2+2x+3$, $x^2-2$ and $x^2+x+1$. Give exact roots, then approximations explicitly marked with $\approx$. Check your results by expanding a factorisation or substituting the roots.\par
\end{document}
